\section{Records and maps}
\label{sec:references}

Y has \emph{two} key/value types, deliberately distinct:

\begin{center}
\begin{tabular}{lll}
\toprule
 & \textbf{record} (static) & \textbf{map} (dynamic) \\
\midrule
type & \ic{\{ a: long, b: str \}} & \ic{\{str: long\}} \\
keys & identifiers, fixed at compile time & any hashable value at run time \\
values & may differ per field & one homogeneous type \ty{V} \\
layout & C struct: fields at fixed offsets & growable key/value store \\
keying & \ic{r.a} & \ic{m['a']} \\
\bottomrule
\end{tabular}
\end{center}

\subsection{The value model}

\begin{center}
\begin{tabular}{lll}
\toprule
category & types & assignment \\
\midrule
values & \ic{long} \ic{byte} \ic{bool} \ic{float} \ic{str} \ic{null}
   & copies the bits \\
references & arrays (\ic{[T]}, \ic{[T; N]}), records, maps
   & \textbf{aliases} the object \\
\bottomrule
\end{tabular}
\end{center}

\noindent ``Reference'' means: every binding of the same object names
one storage, and every in-place mutation (\ic{+=}, \ic{*=}, field
assignment, map stores) is observable through all of them.

\subsection{Records}

A record literal \ic{\{ name: "y", n: 1 \}} allocates one record object
whose shape is the literal's fields \emph{as written}; field order is
preserved. At run time a record is exactly a C struct: field \ty{i} is
the word at \ic{+8i} from the record address, so \ic{r.a} compiles to a
single load at a compile-time offset.

\begin{lstlisting}
let r = { v: [7], n: 1 };
let s = r;            // aliases the record
r.v += 8;             // s.v sees it
s.n = 2;              // r.n sees it
\end{lstlisting}

\noindent The shape is fixed: writing a field the shape does not have is
a \emph{compile error} (there is no runtime field growth), and a field
read requires the field to exist statically.

\paragraph{Record spread.} \ic{\{a: 0, ...obj\}} builds a record whose
shape is the \emph{union} of the two shapes — entirely statically; no
dynamic shape is involved. Entries merge with last-wins values and
first-appearance positions, so \ic{obj} may overwrite \ic{a} and a later
explicit key may overwrite what \ic{obj} spread:

\begin{lstlisting}
let b = {x: 1, y: 2};
let c = {y: 9, z: 3};
print({a: 0, ...b});     // {a: 0, x: 1, y: 2}
print({...b, ...c});     // {x: 1, y: 9, z: 3}
print({...b, y: 100});   // {x: 1, y: 100}
\end{lstlisting}

\subsection{Maps}
\label{sec:maps}

A map type is \ic{\{\ty{K}: \ty{V}\}} where \ty{K} is a hashable key
type (\ic{str}, \ic{long}, \ic{byte}, \ic{bool}, \ic{float}) and \ty{V}
is one homogeneous value type. Map literals use quoted keys (identifier
keys are record syntax):

\begin{lstlisting}
let m = {"a": 1, "b": 2};      // {str: long}
let empty: {str: long} = {};   // empty literals need the annotation
m["c"] = 3;                    // insert or overwrite
print(m["a"]);                 // 1
print(m["zz"]);                // null -- reads yield V?
\end{lstlisting}

\noindent Keying is symbolic and exclusive: \ic{m['k']} is the map form
and is reserved for maps; \ic{r.a} is the record form. Reading
\ic{m[k]} yields \ic{\ty{V}?} — \ic{null} when the key is absent — and
\ic{m[k] = v} stores in place (the store is visible through every
binding of the map). \ic{len(m)} is the number of entries. The map grows
as needed (amortized doubling, like arrays).

\subsection{Equality}

\ic{==} on two arrays or two records is \textbf{structural}: same
length/fields, pairwise-equal contents. It is never a pointer
comparison; \ic{arr == []} is an emptiness test.

\subsection{Copying}

There is no implicit deep copy. To duplicate an array use \ic{+} with an
empty right side (\ic{a + []}) or slicing (\ic{a[:]}); records can be
rebuilt field by field. The language deliberately makes sharing the
default and copying explicit.
